\section{Introduction}
When each image has a fixed channel-symbol budget, the transmitter must decide which source information to retain, and the receiver must reconstruct the image from that description. Feasibility depends on the image-dependent source-code length and the protection carried within the frame. A generative prior can complete untransmitted information, while the delivered representation provides source-image constraints.

Digital coding with error correction and learned joint source--channel coding (JSCC) provide different solutions~\cite{deepjscc,swinjscc}; fitting progressive descriptions to channel resources is a classical problem~\cite{chande}. Ada-TokenCom combines prefix entropy coding, tail generation, and prefix/MCS selection under an average resource constraint~\cite{adatokencom}. TONIC studies task-oriented token communication~\cite{tonic}, while ARPC constructs progressive descriptions with an autoregressive prior~\cite{arpc}. We adopt the next-scale representation of visual autoregressive modeling (VAR)~\cite{var}. Complete coarse maps provide a low-resolution description of the full latent grid; subsequent maps add residual information and are predicted from earlier scales. This organization supports reconstruction from a coarse description when only a limited prefix can be delivered.

Coarse-to-fine refinement does not make whole-scale boundaries align with digital-frame capacity. Adding a complete scale requires a large token increment, and the independently terminated entropy-code length of a given endpoint varies across images. After one complete scale fits, the remaining capacity may be insufficient for the next. Within-scale endpoints retain source tokens from that boundary map. The next-scale factorization lets both ends construct matching coding probabilities from complete coarse maps. To turn the additional received information into reconstruction gains, the boundary tokens must enter residual accumulation and the conditions used to predict finer scales.

NeST-Com combines within-scale endpoints, independently terminated arithmetic streams, actual-length fitting, and prefix-constrained reconstruction in a fixed-budget digital system. It uses VAR's known-token overwrite operation~\cite{var} to merge received boundary positions with predictions and update both the residual latent and finer-scale context. With shared model components, PSNR gains over whole-scale transmission are 0.32--0.70~dB at $N=1024$ and 0.39--0.79~dB at $N=512$; utilization rises from 77.70\% to 95.21\% at $N=1024$, 10~dB. Controlled endpoint and receiver interventions examine these design gains. Six-system N1024 comparisons and independently calibrated N512 results show the source-reconstruction performance and perceptual tradeoffs of the complete systems.
